
%%% Local Variables:
%%% mode: latex
%%% TeX-master: t
%%% End:

\chapter{实验原理与模型构建}
\label{cha:command}

\section{实验原理}
\subsection{前馈神经网络}
深度神经网络是一种多层前馈神经网络，其基本结构包括输入层、隐藏层和输出层，每一层的计算形式为
\begin{equation*}
    y_{k+1}=\sigma(w^T x_k+b),
\end{equation*}
其中$w$为权重矩阵，$b$为偏执，$\sigma$为激活函数.

\subsection{激活函数}
激活函数用于引入非线性能力. 本实验研究三种常见激活函数：ReLU、tanh与sigmoid

\subsection{损失函数}
本实验使用均方误差（MSE）：
\begin{equation*}
    \mathbf{MSE} = \frac{1}{n}\sum_{i=1}^n(y_i-\hat{y_i})^2,
\end{equation*}
并计算平均绝对误差(MAE)：
\begin{equation*}
    \mathbf{MAE} = \frac{1}{n}\sum_{i=1}^N\left | y_i-\hat{y_i}\right |.
\end{equation*}

\subsection{优化算法}
本实验使用Adam优化器



\section{模型构建}
\subsection{实验环境与版本配置}
实验环境如下
\begin{table}[htbp]
  \centering
  \caption{表格示例}
    \begin{tabular}{ccccc}
    \toprule
    项目 & 配置 \\
    \midrule
    Python版本  & 3.14.0 \\
    深度学习框架  & PyTorch \\
    numpy  & 2.4.3 \\
    pandas  & 2.3.3 \\
    scipy & 1.17.1 \\
    matplotlib & 3.10.8 \\
    \bottomrule
    \end{tabular}%
  \label{Tab4-1}%
\end{table}%

\subsection{模型配置}
本实验默认配置如下表：
\begin{table}[htbp]
  \centering
  \caption{模型默认参数}
    \begin{tabular}{ccccc}
    \toprule
    项目 & 配置 \\
    \midrule
      网络结构 & [1,50,50,1] \\
      激活函数 & tanh \\
      训练回合数 & 30000 \\
      总样本点数 & 151 \\
      训练点数 & 75 \\
    \bottomrule
    \end{tabular}%
  \label{Tab4-1}%
\end{table}%

根据实验内容，变更相应实验参数，其余保持不变